\section{Related Work}

\subsection{Most Directly Related: Dynamic and Graph-based ANNS}
Graph-based methods such as HNSW and its variants are widely adopted for high-recall and low-latency vector search.
Most designs, however, are optimized for static or slowly changing datasets, where index quality and auxiliary search artifacts are built offline.
Recent dynamic ANNS studies explore incremental index maintenance and update-aware operations.
In contrast, StreamSeed focuses on policy-level control of query hint reuse and provides a pluginized path to integrate with existing dynamic backends.

\subsection{Indirectly Related: System and Algorithm Acceleration}
Prior system works improve ANNS throughput via hardware-conscious layouts, vector quantization, and SIMD-friendly distance computation.
These approaches are complementary to StreamSeed: they optimize execution efficiency but do not directly address online staleness of reusable query guidance.
Algorithmic methods reduce unnecessary distance evaluations through pruning, lower-bound estimation, or partial-dimension strategies.
While effective in static contexts, they usually rely on assumptions that weaken under high-rate updates and distribution drift.

\subsection{Other: Positioning of This Work}
StreamSeed is orthogonal to both system-only and algorithm-only accelerations.
Its contribution is a unified online control loop (AQHC) that jointly manages hint construction, usage, and update,
bridging algorithmic reuse opportunities with systems-level deployability in streaming settings.
